\documentclass[11pt]{article}

\usepackage[margin=1in]{geometry}
\usepackage[T1]{fontenc}
\usepackage[utf8]{inputenc}
\usepackage{lmodern}
\usepackage{microtype}
\usepackage{xcolor}
\usepackage{hyperref}
\usepackage{listings}
\usepackage{setspace}

\hypersetup{
  colorlinks=true,
  linkcolor=blue!50!black,
  citecolor=blue!50!black,
  urlcolor=blue!60!black,
  pdftitle={Register-Token Sidecar for Efficient Shadow-Mode Data Pipelines in End-to-End Autonomous Driving},
  pdfauthor={[YOUR FULL NAME HERE]}
}

\setlength{\parindent}{0pt}
\setlength{\parskip}{0.65em}
\setstretch{1.05}

\makeatletter
\renewcommand\@seccntformat[1]{\csname the#1\endcsname.\quad}
\makeatother

\pagestyle{plain}

\lstdefinestyle{sidecarcode}{
  language=Python,
  basicstyle=\ttfamily\small,
  keywordstyle=\color{blue!60!black}\bfseries,
  commentstyle=\color{teal!50!black}\itshape,
  stringstyle=\color{orange!50!black},
  showstringspaces=false,
  frame=single,
  rulecolor=\color{black!20},
  backgroundcolor=\color{black!2},
  xleftmargin=0.5em,
  xrightmargin=0.5em,
  aboveskip=1em,
  belowskip=1em,
  breaklines=true,
  columns=fullflexible,
  numbers=left,
  numberstyle=\tiny\color{black!45},
  numbersep=8pt
}

\title{\textbf{Register-Token Sidecar for Efficient Shadow-Mode Data Pipelines in End-to-End Autonomous Driving}}
\author{[YOUR FULL NAME HERE]}
\date{May 1, 2026}

\begin{document}

\maketitle
\thispagestyle{empty}

\begin{center}
\small
\textit{Theory \& Validation Blueprint}

\textit{Independent research proposal. Not affiliated with Tesla, Valeo, or the authors of DrivoR.}
\end{center}

\begin{abstract}
End-to-end autonomous driving systems rely on massive fleets for continuous improvement through shadow-mode data collection. While current trigger heuristics effectively reduce upload volume compared to raw video, two open questions remain: (1) whether an orthogonal signal can surface long-tail failures that existing triggers miss, and (2) whether uploaded payloads can be further compressed while preserving downstream training signal.

We propose a lightweight register-token sidecar inspired by the DrivoR architecture \cite{kirby2026drivor}. A LoRA-tuned head with 16--32 camera-aware register tokens is run in inference-only mode alongside the main policy network. We formalize two testable hypotheses: register statistics as a hard-example mining signal (H1) and register tokens as a semantic compression codec (H2). We provide a complete 3-week offline validation plan, pseudocode, clear kill criteria, and risk assessment. No changes are made to the core pixels-to-controls policy network.

This sidecar is designed to be a low-risk, high-leverage data-pipeline upgrade compatible with any large-scale end-to-end fleet.
\end{abstract}

\section{Introduction}

Modern end-to-end autonomous driving programs, including fleet-scale systems such as Tesla FSD v12+ and similar stacks, have largely eliminated explicit maps and symbolic world models in favor of pure neural policies trained on vast fleet data. Shadow-mode collection --- uploading only ``interesting'' clips triggered by disengagements, planner disagreement, or heuristics --- is now standard practice to manage bandwidth and storage.

However, at planetary fleet scale, even highly optimized triggers leave room for improvement in two dimensions:

\begin{enumerate}
  \item \textbf{Orthogonality of hard-example discovery} --- Some long-tail failures remain under-sampled.
  \item \textbf{Payload efficiency} --- The current format (triggered imagery + dense features) is still expensive.
\end{enumerate}

This whitepaper proposes a minimal, non-invasive sidecar that addresses both issues using the recently published DrivoR register-token primitive.

Concretely, the proposal contributes three things: a low-risk sidecar integration point, two testable hypotheses for data-pipeline improvement, and a short offline-first validation plan with explicit decision gates.

\section{Background: DrivoR Register Tokens}

DrivoR \cite{kirby2026drivor} introduces camera-aware register tokens into a pretrained DINOv2 ViT-S backbone. A small LoRA head (1--3 million trainable parameters) is added to produce 16--32 compact per-camera registers that replace the full set of patch tokens.

Key published results:

\begin{itemize}
  \item More than 250x token compression
  \item No loss in trajectory prediction quality on NAVSIM v1/v2 and closed-loop HUGSIM
  \item Forward pass approximately 110 ms on a single A100 (single batch)
\end{itemize}

The registers act as a learned, compact scene descriptor while remaining fully compatible with downstream transformer decoders. We repurpose only the register mechanism as a passive sidecar for data efficiency.

\section{Problem Statement}

Current shadow-mode pipelines are effective but not exhaustive. We ask:

\textbf{Question 1 (Orthogonality):} Do register-token statistics provide a signal that correlates with real interventions above and beyond existing triggers?

\textbf{Question 2 (Compression):} Can register tokens + minimal triggered imagery serve as a drop-in replacement payload that preserves equivalent training signal?

\section{Proposed Approach}

We attach a frozen DINOv2 ViT-S + LoRA register head to the existing perception stack. The head runs in parallel in shadow mode only.

\begin{description}
  \item[\textbf{H1 --- Hard-Example Mining}] Compute scalar statistics on the register tokens (assignment entropy, inter-register variance, distance to learned prototypes, etc.). High-score clips are hypothesized to be enriched for interventions missed by existing triggers.
  \item[\textbf{H2 --- Semantic Compression}] Transmit only the compact register tokens (a few KB) plus minimal triggered frames.
\end{description}

\section{Experimental Validation Plan}

\begin{table}[htbp]
\centering
\small
\begin{tabular}{|p{0.22\textwidth}|p{0.18\textwidth}|p{0.5\textwidth}|}
\hline
\textbf{Phase} & \textbf{Window} & \textbf{Primary Objective} \\
\hline
Offline validation & Weeks 0--3 & Test orthogonality and payload efficiency using public datasets or archived fleet clips. \\
\hline
Shadow-mode A/B & Weeks 4--7 & Conditional deployment on a limited fleet if offline metrics justify progression. \\
\hline
Decision gate & End of Phase 2 & Continue only if upload reduction, latency, and orthogonality targets are all met. \\
\hline
\end{tabular}
\end{table}

\textbf{Phase 1 --- Offline Validation (Weeks 0--3)}

Use public datasets (nuScenes, BDD100K) or archived fleet clips.

\begin{itemize}
  \item Run DrivoR register head on 1--10 million frames.
  \item H1: ROC-AUC of register-derived scores vs.\ intervention labels, conditioned on existing-trigger = negative.
  \item H2: Train lightweight downstream head on register-only vs.\ baseline payloads.
  \item Measure forward-pass latency on HW4-class hardware.
\end{itemize}

\textbf{Phase 2 --- Shadow-Mode A/B (Weeks 4--7, conditional)}

Deploy inference-only sidecar on 50--100k vehicles.

\textbf{Decision Gate (Kill Criteria)}

Proceed only if all hold:

\begin{itemize}
  \item At least 3x reduction in upload volume at parity intervention-discovery rate
  \item Added latency within Autopilot budget (less than 3 ms target on HW4)
  \item Statistically meaningful orthogonality (H1)
\end{itemize}

\section{Illustrative Pseudocode}

\begin{lstlisting}[style=sidecarcode, caption={Inference-only register-token sidecar workflow.}, label={lst:sidecar}]
vit = load_dinov2_s_backbone(frozen=True)
register_head = LoRA_RegisterHead(num_registers=16)  # DrivoR weights

for clip in dataset:
    features = vit(clip.frames)
    registers = register_head(features)

    entropy = compute_assignment_entropy(registers)
    variance = compute_inter_register_variance(registers)
    proto_dist = compute_distance_to_prototypes(registers)

    score = combine_scores(entropy, variance, proto_dist)
    log(score, clip.intervention_label, clip.existing_trigger_label)
\end{lstlisting}

\section{Risks and Mitigations}

\begin{itemize}
  \item \textbf{Latency on HW4} --- measured first.
  \item \textbf{Low orthogonality} --- primary H1 kill condition.
  \item \textbf{Training-signal degradation (H2)} --- verified with downstream head.
  \item \textbf{Scope creep} --- strictly a data-pipeline sidecar.
\end{itemize}

\section{Conclusion}

The register-token sidecar offers a low-risk, high-leverage upgrade for any end-to-end autonomous driving fleet. Planned deliverables include open-source code and checkpoints upon completion of the offline phase.

We invite feedback, collaboration, or joint validation.

\begin{thebibliography}{9}

\bibitem{kirby2026drivor}
Kirby, E., et al.
\newblock \emph{Driving on Registers}.
\newblock arXiv:2601.05083, Valeo.ai, January 2026.
\newblock Presented at CVPR 2026.
\newblock Available at: \url{https://arxiv.org/abs/2601.05083}.

\end{thebibliography}

\end{document}
